\section{Introduction}\label{sec:intro}

\subsection{Motivation}\label{sec:intro:motivation}

Autonomous research systems promise to shorten the path from a question
to an experimentally grounded answer. Their central difficulty is not
merely generating code or prose. A useful system must decide what to try,
judge noisy outcomes, preserve evidence, recover from failure, and keep a
manuscript from claiming more than the run established. These decisions
form an institution: proposers, evaluators, adversaries, writers, and
judges whose behaviour is largely encoded in prompts and scoring rules.

A fixed institution is reproducible but inherits every weakness of its
initial design. An evolvable institution can adapt, but creates a reflexive
control problem. A component that benefits from a changed score may supply
the evidence for changing it; a reviewer may grade its own successor; a
prompt that gains tool access may hide the history that exposes the
expansion. Self-improvement is therefore an authority problem as well as
an optimisation problem.

\subsection{Verified claims as a release condition}\label{sec:intro:verified}

Constitutional Autonomous Research Infrastructure (ARI)-RQGM starts from a
narrow rule: model outputs may
propose institutional change, but no model output may enact it directly.
A deterministic kernel validates record shapes, hashes, capabilities,
role separation, transition legality, audit continuity, clean-room
inputs, and context scope. A single transition engine owns registry
writes. Changes occur only at boundaries; T16 emergency quarantine
force-closes the current epoch and opens a newly fingerprinted epoch in the
same transaction. A boundary transaction either commits in full or
disappears on replay.

The fixed layer also retains the final structured numeric
claim--evidence decision. A
governed writer and reviewer may compare manuscript alternatives, but
neither can modify the verifier that recomputes forward-declared formulas
from executed evidence. It is not a general natural-language truth
checker. Governance selects what is presented to the gate; it cannot
relax the gate.

\subsection{Contributions}\label{sec:intro:contrib}

This report makes five contributions.

\begin{enumerate}
  \item It defines a three-layer authority model: an immutable
        constitutional layer, an evolvable institutional layer, and a
        meta layer that may propose replacements but cannot appoint them.
  \item It specifies the complete T1--T21 component-lifecycle table,
        including the forced emergency boundary for T16, together with the capability
        matrix, event contract, and deterministic recovery procedure.
  \item It closes the accountability loop from a raw artefact attack to
        defence, adjudication, a validated attack, and target binding for
        the epoch-stamped research generator and the currently connected
        paper roles. Ambiguous provenance remains deliberately targetless.
  \item It gives retirement operational meaning through selective
        erasure, clean-room regeneration, replay evaluation, and
        provenance-aware repair, including successor-policy re-scoring
        from stored raw axes.
  \item It applies the constitution to a governed paper archive and
        reports lifecycle and fault-injection evidence from the
        implementation.
\end{enumerate}

\subsection{Scope of this report}\label{sec:intro:scope}

The subject is institutional co-evolution, not a catalogue of every
execution feature in ARI. We treat the experiment executor, domain tools,
and compilation pipeline as an environment that produces and consumes
typed artefacts. Their internal scheduling strategies are out of scope.
The constitutional claims should hold for any research executor that
obeys the record, provenance, and checkpoint contracts.

We distinguish construction guarantees from empirical claims.
Single-writer registry authority, legal transition topology, and
fail-closed state changes are structural. A local hash chain detects
interior modification but not rollback of the whole checkpoint and local
pin to an old valid prefix. Improved
novelty, research quality, calibration, and cost effectiveness require
controlled evidence. Current results are 1,124 passing directly relevant
tests, five fixed fault fixtures, and one clean fixture. They do not
establish general detection rates, attribution accuracy, research-quality
improvement, or cost effectiveness.

\subsection{Organisation}\label{sec:intro:organisation}

\Cref{sec:overview} defines the system and threat model.
\Cref{sec:rqgm} specifies the constitutional lifecycle, co-evolution,
paper governance, and observability. \Cref{chapter:eval} separates
demonstrated behaviour from deferred claims. \Cref{sec:related} positions
the architecture against self-modifying agents, autonomous research
systems, and verification work. \Cref{sec:discussion} discusses the
fixed-constitution trade-off and open scaling questions. The appendices
give the complete core specification and reproduce the RQGM-specific
runtime prompt seeds verbatim.
